Realizace experimentů předpokládá, že máme k dispozici materiál: grafen oxid, předmět měření: membránu, měřicí aparaturu: permeametr, a v neposlední řadě záznamový i vyhodnocovací software, kde míra překryvu jednotlivých prvků sekvence se specializací autora ve stejném pořadí roste. Pro materiálového inženýra či fyzikálního chemika tak text Experimentální části může místy působit jako přehnaně, či naopak nedostatečně, detailní; jde o nevyhnutelný důsledek širokého mezioborového přesahu práce.

Kapitola \ref{s:grafen_oxidova_membrana} popisuje, po vzoru materiálové části vědeckých publikací, přípravu grafen oxidu a výrobu membrány. Pro další měření použitý grafen oxid pocházel ze tří šarží, z nichž chronologicky prostřední je syntézou autora. Veškerý použitý grafen oxid byl syntetizován Tourovou, zvanou též Vylepšenou Hummersovou, metodou; k výrobě membrán byla použita výhradně přetlaková filtrace. Syntéza všech šarží byla realizována v laboratořích prof. Sofera, k čištění a charakterizaci materiálu docházelo dle aktuálních kapacit též v Laboratoři membránových separačních procesů prof. Friesse. 

V Kapitole \ref{s:integralni_manometricky_permeametr} se věnujeme vývoji měřicí aparatury. První funkční prototyp integrálního manometrického permeametru s řízeným vstupním tlakem byl v Laboratoři membránových a separačních procesů postaven zhruba v roce 2017. S~ohledem na vývoj požadavků měření, zejména větší zaměření se na vliv vstupního tlaku, byla aparatura v několika etapách modifikována až do prezentované podoby. Z~původní aparatury zůstávají základní stavební prvky, jako jsou samotné tělo cely, tlakoměr či regulátor; na druhou stranu fyzické uspořádání, řešení těsnění či otázka termostatování aparaturu zásadně pozměnily.

V Kapitole \ref{s:software} popisujeme vývoj programu, který permeametr ovládá. Program k původnímu permeametru byl psán v prostředí LabVIEW, fungoval pouze jako zařízení pro zápis dat ze senzoru; řízení vstupního tlaku probíhalo externě na dedikovaném zařízení, k zápisu dat z regulátoru nedocházelo vůbec. Vyvinutý program propojuje v jednom prostředí ovládání regulátoru i senzoru, poskytuje různé módy měření, nabízí online vizualizaci a okamžité vyhodnocení měření; nadto realizuje asynchronní zápis dat a oproti původnímu zásadně zrychuje vzorkovací frekvenci. Kód je od vlastní implementace ovladačů až po grafické rozhraní autorským řešením.

